\subsection{Prediction 3: Why Japan Lost Decades While Sweden Lost Only Years}
\label{subsec:pred3-asymmetric}

%%%%%%%%%%%%%%%%%%%%%%%%%%%%%%%%%%%%%%%%%%%%%%%%%%%%%%%%%%%%%%%%%%%%%%%%%%%%%%%
% CHAPTER 11.3: Narrative Treatment of Prediction 3
% Intuitive explanation of convex recovery and the big push principle
%%%%%%%%%%%%%%%%%%%%%%%%%%%%%%%%%%%%%%%%%%%%%%%%%%%%%%%%%%%%%%%%%%%%%%%%%%%%%%%

\subsubsection{Two Countries, One Year, Two Fates}

In 1991, two prosperous northern countries suffered banking crises.

In Japan, a massive real estate bubble burst. Stock prices crashed. Banks 
found themselves holding worthless loans. The government responded with 
stimulus packages, interest rate cuts, public works programs.

In Sweden, a similar pattern: property bubble, banking collapse, economic 
contraction. The government nationalized failing banks, created a ``bad bank'' 
to absorb toxic assets, and recapitalized the financial system.

Japan's crisis was bigger---about 30\% larger by most measures. Sweden's 
response was faster and more aggressive. But the difference in outcomes 
was far larger than these differences would suggest.

Sweden recovered in 5--7 years. By 1998, it was back on track.

Japan lost not one decade, but two. As of 2024---thirty-three years later---Japan 
has never fully returned to its pre-crisis growth path.

How can a 30\% larger crisis produce a 300\% longer recovery?

\subsubsection{The Convexity Insight}

The answer lies in a single word: \textbf{convexity}.

Standard economic thinking assumes proportional relationships. Twice the 
shock, twice the recovery time. This linearity is built into most policy 
models, and it profoundly shapes crisis response.

But recovery isn't linear. It's convex.

Small crises recover quickly because the economic machinery is still mostly 
intact. Banks still trust each other. Workers still have skills. Supply 
chains still function. Confidence, though shaken, remains. The economy 
can heal itself with modest support.

Large crises are fundamentally different. Below a certain threshold, the 
machinery breaks. Banks don't just reduce lending---they stop entirely. 
Workers don't just change jobs---they drop out of the workforce. Confidence 
doesn't just dip---it collapses into self-fulfilling pessimism.

Rebuilding broken machinery takes far longer than healing wounded machinery.

\subsubsection{The Valley of Incoherence}

We can visualize this with what we call the ``coherence valley'':

\begin{center}
\textit{[Imagine a valley with steep sides]}
\end{center}

\begin{itemize}
    \item At the top of the slope, the economy is functioning. Small pushes 
    knock it down a bit, but it naturally rolls back up.
    
    \item Partway down, the slope gets steeper. The economy can still recover, 
    but it takes more effort.
    
    \item At the bottom is a flat valley floor. Once the economy falls here, 
    it doesn't naturally roll anywhere. It just sits there.
    
    \item Getting out of the valley requires climbing the steep slope---much 
    harder than the initial fall down.
\end{itemize}

Sweden fell partway down the slope. Gravity (natural market forces) plus 
a good push (policy intervention) was enough to get back up.

Japan fell into the valley. No amount of gradual pushing could climb 
that slope. Only a massive, coordinated effort---which never came---could 
have lifted it out.

\subsubsection{The Numbers}

Our framework quantifies this convexity. We define ``coherence'' ($K$) as 
a measure of how well the economic system's parts work together. $K = 1$ 
means everything is functioning smoothly. $K = 0$ means complete breakdown.

The critical insight: recovery time isn't proportional to how far $K$ falls. 
It's proportional to something steeper:

\begin{center}
\renewcommand{\arraystretch}{1.3}
\begin{tabular}{lccc}
\hline
\textbf{Crisis Severity} & \textbf{Post-Crisis $K$} & \textbf{Recovery Time} & \textbf{vs. Mild Crisis} \\
\hline
Mild & 0.85 & $\sim$2 years & 1× \\
Moderate & 0.70 & $\sim$5 years & 2.3× \\
Severe & 0.55 & $\sim$9 years & 4.4× \\
Very Severe & 0.40 & $\sim$19 years & 8.8× \\
\hline
\end{tabular}
\end{center}

A very severe crisis (dropping $K$ to 0.40) doesn't take twice as long 
as a moderate crisis (0.70). It takes \emph{four times} as long.

This is convexity in action.

\subsubsection{The Threshold That Changes Everything}

There's another critical feature: a threshold around $K \approx 0.50$.

Above this threshold, the economy has enough functioning components that 
normal policy works. Cut interest rates, stimulate demand, clean up bad 
loans---and recovery follows.

Below this threshold, normal policy stops working. The complementarities 
that make an economy function---bank trust, consumer confidence, business 
investment, institutional credibility---have broken down. Fixing one doesn't 
help because the others are still broken.

Think of it like a car with multiple problems. If the battery is weak, 
jump-starting it works. But if the battery is dead, the starter motor is 
burned out, and the fuel line is clogged---jump-starting does nothing. 
You need to fix everything at once.

Our estimates suggest:
\begin{itemize}
    \item Sweden post-crisis: $K \approx 0.58$ (above threshold)
    \item Japan post-crisis: $K \approx 0.42$ (below threshold)
\end{itemize}

Sweden was damaged but recoverable with standard tools.

Japan had fallen below the threshold where standard tools work.

\subsubsection{The Big Push Principle}

When you're below the threshold, only a ``big push'' works.

The term comes from development economics, where it describes the idea 
that poor countries can't develop gradually---they need simultaneous 
investment in infrastructure, education, institutions, and industry, 
because each depends on the others.

The same logic applies to crisis recovery. Below $\underline{K}$:

\begin{enumerate}
    \item \textbf{Marginal reforms produce marginal results.} Cleaning up 
    10\% of bad loans doesn't restore 10\% of lending---it restores almost nothing, 
    because the system remains broken.
    
    \item \textbf{Coordination is essential.} Banks, firms, consumers, and 
    government must all act together. If each waits for the others, nothing happens.
    
    \item \textbf{Expectations must shift.} People must believe recovery is 
    coming for their behavior to make recovery come. Half-measures don't 
    change expectations.
\end{enumerate}

Sweden delivered something close to a big push: nationalization, bad bank, 
recapitalization, all at once. It was expensive and controversial, but 
it worked.

Japan delivered incremental responses: stimulus packages, partial cleanups, 
interest rate cuts. Each was reasonable in isolation. Together, they were 
inadequate.

\subsubsection{Why Japan Didn't Push Big}

This raises an obvious question: if big push was needed, why didn't Japan do it?

The answer reveals the cruel political economy of crisis response.

\paragraph{1. They didn't know they were below threshold.}
In 1991, no one knew this would become the ``lost decades.'' Officials 
believed the crisis was serious but manageable---a $K \approx 0.65$ situation 
requiring $K \approx 0.65$ responses. By the time the true severity was clear, 
$K$ had fallen further.

\paragraph{2. Big push is politically hard.}
A big push is expensive, visible, and concentrated. The costs appear 
immediately; the benefits take years. Politicians facing elections have 
strong incentives to prefer gradual, visible action (``we're doing something!'') 
over dramatic intervention that might not pay off until after they're out of office.

\paragraph{3. Institutions constrained action.}
Japan's consensus-based political system made rapid, comprehensive action 
difficult. Coordinating across ministries, parties, and interest groups 
took time---time during which $K$ kept falling.

\paragraph{4. Zombie firms complicated everything.}
Banks had incentives to keep failing firms alive (to avoid recognizing losses). 
This ``zombie lending'' prevented the creative destruction that recovery required. 
But stopping it meant acknowledging how bad things were.

Sweden, by contrast, had:
\begin{itemize}
    \item Clearer recognition of severity (earlier, accurate diagnosis)
    \item Stronger state capacity (faster implementation)
    \item More political consensus (easier coordination)
    \item Willingness to let failing firms fail
\end{itemize}

\subsubsection{The Welfare Cost of Getting It Wrong}

The stakes of this asymmetry are enormous.

We can calculate the cumulative welfare loss from Japan's slow recovery 
versus the counterfactual of Swedish-speed recovery:

\begin{itemize}
    \item Japan's actual growth (1991--2020): $\sim$1.0\% per year
    \item Counterfactual growth (Swedish-speed recovery): $\sim$2.5\% per year
    \item Gap: $\sim$1.5\% per year for 30 years
    \item Cumulative: $\sim$45\% of 1991 GDP that never materialized
\end{itemize}

In human terms: a generation of Japanese workers earned less, saved less, 
and retired with less than they would have with faster recovery. A generation 
of young people entered a stagnant job market and never fully recovered. 
Birth rates fell. Pessimism became normalized.

All because the initial response was calibrated to a moderate crisis, 
not the severe crisis that was actually happening.

\subsubsection{The Lesson for Future Crises}

This analysis leads to a clear but uncomfortable prescription:

\begin{tcolorbox}[colback=gray!10!white, colframe=gray!75!black]
\textbf{When in doubt, overshoot.}

The cost of undershooting a severe crisis is decades of stagnation.
The cost of overshooting a moderate crisis is some wasted spending.

These are not symmetric risks.
\end{tcolorbox}

The 2008 global financial crisis tested this lesson. The US response 
was larger and faster than typical; the Eurozone response was smaller 
and slower. The US recovered in $\sim$5 years; parts of Europe took 10+.

COVID-19 (2020) tested it again. Countries that responded with overwhelming 
force (massive stimulus, comprehensive support) recovered faster than those 
that responded incrementally.

The pattern keeps confirming: in severe crises, go big or go home.

\subsubsection{Summary: Prediction 3}

\begin{tcolorbox}[colback=gray!5!white, colframe=gray!75!black, title=Prediction 3: Asymmetric Recovery]
\textbf{What standard theory predicts:}
\begin{itemize}
    \item Recovery time is proportional to crisis magnitude
    \item Policy works similarly in mild and severe crises
    \item Gradual adjustment is generally optimal
\end{itemize}

\textbf{What the framework predicts:}
\begin{itemize}
    \item Recovery is convex: big crises take disproportionately longer
    \item Below threshold $\underline{K} \approx 0.50$, normal policy fails
    \item Only ``big push'' coordinated intervention works for severe crises
    \item Undershooting is far costlier than overshooting
\end{itemize}

\textbf{What the data show:}
\begin{itemize}
    \item Japan vs. Sweden: 30\% larger crisis → 300\% longer recovery ✓
    \item Convexity coefficient: $\beta_2 = 0.45 > 0$ (significant) ✓
    \item Threshold exists: $\underline{K} \approx 0.52$ (estimated) ✓
    \item Pattern holds across 70 crises since 1970 ✓
\end{itemize}

\textbf{Bottom line:} Big crises need disproportionately big responses. 
Japan's lost decades weren't inevitable---they resulted from treating a 
severe crisis with moderate tools. When you're below the threshold, 
only a big push works.
\end{tcolorbox}

\vspace{0.5em}
\noindent\textit{Technical details: See Appendix AWA, Section~\ref{subsec:pred-03-calc} for the 
formal convexity model, threshold estimation, Japan-Sweden comparison, and 
policy calibration framework.}

%%%%%%%%%%%%%%%%%%%%%%%%%%%%%%%%%%%%%%%%%%%%%%%%%%%%%%%%%%%%%%%%%%%%%%%%%%%%%%%
